
\section{Aufgabe 2}

Die Messungen wurden auf einer NVIDIA GeForce RTX 4060 durchgeführt. Laut
Geräteabfrage liegt die theoretische Speicherbandbreite bei
272 GB/s. Die GPU besitzt 24 SMs und maximal 1536 Threads pro SM, also
48 Warps pro SM.

\subsection{Speicherintensiver Kernel}

Ein speicherintensiver Kernel wurde in \texttt{memory.cu} implementiert. Jeder Thread 
liest einen Wert aus \texttt{a} multipliziert diesen mit einer Konstanten \texttt{c}
und schreibt diesen Wert nach \texttt{b}. 

\begin{verbatim}
b[i] = a[i] * c;
\end{verbatim}

Während der Rechenaufwand pro Element mit höchstens 0{,}25 FLOP/Byte sehr gering ist, 
entstehen Speicherzugriffe im Umfang von 8 Byte durch einen 4 Byte \textit{read} und 
einen 4 Byte \textit{write}. Damit wird die Laufzeit des Kernels primär durch das 
Speicherzugriffsmuster und nicht die ALUs bestimmt.


Es wurden insgesamt drei Varianten des Kernels implementiert:

\begin{itemize}
  \item \texttt{kernel\_copy\_coalesced}: benachbarte Threads lesen
        benachbarte Adressen.
  \item \texttt{kernel\_copy\_strided}: Thread \texttt{i} liest
        \texttt{a[(i * stride) \% n]}.
  \item \texttt{kernel\_copy\_gather}: Thread \texttt{i} liest eine
        zufällige Adresse \texttt{a[idx[i]]}.
\end{itemize}

\subsubsection{STREAM-Triad als Referenzfall}

Zusätzlich wurde ein STREAM-Kernel in \texttt{stream.cu} als Referenz implementiert. 
Dieser bearbeitet ein klassisches Bandbreitenproblem:
%Der STREAM-Kernel ist deshalb ein sehr guter Referenzfall
%für ideale, zusammenhängende und coalesced Speicherbandbreite.

\begin{verbatim}
a[i] = b[i] + alpha * c[i];
\end{verbatim}

Pro Element werden zwei Werte \texttt{b[i]} und \texttt{c[i]} gelesen und ein Wert 
\texttt{a[i]} geschrieben. Das erzeugt bei FP32 Speicherzugriffe im Umfang von 12 Bytes 
für nur eine \textit{fused multiply add (fma)}. Diese hat einen arithmetischen Aufwand von  
etwa 0{,}17 FLOP/Byte.

Der Unterschied zu \texttt{memory.cu} besteht darin, dass

\begin{itemize}
  \item \texttt{stream.cu} einen einzelnen, gutartigen
        STREAM-Triad-Fall misst. Bei dem alle Threads regulär und zusammenhängend
        auf \texttt{a}, \texttt{b} und \texttt{c} zugreifen. Der Kernel beantwortet damit die Frage:
        Welche Bandbreite ist bei einem sauberen Streaming-Zugriff erreichbar?
  \item \texttt{memory.cu} variiert dagegen gezielt das Zugriffsmuster. Derselbe
        einfache Rechenkern wird als \textit{coalesced}, \textit{strided} und 
        \textit{random gather} ausgeführt. Damit beantwortet der Kernel die Frage:
        Wie stark bricht die Bandbreite ein, wenn die Adressen innerhalb einer Warps ungünstig liegen?
\end{itemize}

\subsubsection{Zugriffsmuster und effektive Bandbreite}

Für den Sweep wurden $33\,554\,432$ Elemente verwendet. Die
Messwerte liegen in \texttt{patterns.csv}; die daraus erzeugte Grafik ist in
Abbildung~\ref{fig:patterns} dargestellt.

\begin{figure}[H]
  \centering
  \includegraphics[width=0.95\linewidth]{patterns.png}
  \caption{Effektive Bandbreite für coalesced, strided und Gather-Zugriffe.}
  \label{fig:patterns}
\end{figure}

\begin{table}[H]
  \centering
  \begin{tabular}{p{4.2cm}
                  >{\raggedleft\arraybackslash}p{2.5cm}
                  >{\raggedleft\arraybackslash}p{3.0cm}
                  >{\raggedleft\arraybackslash}p{2.5cm}}
    \hline
    \textbf{Muster} & \textbf{Stride} & \textbf{GB/s} & \textbf{\% Peak} \\
    \hline
    coalesced & 1   & 251{,}30 & 92{,}38 \\
    strided   & 2   & 169{,}15 & 62{,}18 \\
    strided   & 4   & 101{,}76 & 37{,}41 \\
    strided   & 8   & 56{,}34  & 20{,}71 \\
    strided   & 16  & 56{,}31  & 20{,}70 \\
    strided   & 32  & 56{,}30  & 20{,}70 \\
    strided   & 64  & 51{,}69  & 19{,}00 \\
    strided   & 128 & 40{,}87  & 15{,}02 \\
    random gather & - & 28{,}03 & 10{,}30 \\
    \hline
  \end{tabular}
  \caption{Gemessene effektive Bandbreite für verschiedene Zugriffsmuster.}
  \label{tab:patterns}
\end{table}

Der \textit{coalesced} Fall (\textit{stride}$=1$) erreicht mit $251{,}30$ GB/s 
etwa $92$ Prozent des theoretischen Peaks. Das ist der Idealfall für eine GPU: 
Alle $32$ Lanes eines Warps lesen zusammenhängende Adressen, sodass die Hardware 
wenige große Speichertransaktionen bilden kann. Damit ist fast jedes übertragene 
Byte nutzbar.

Bei \textit{strided} Zugriffen (\textit{stride}$>1$) sinkt die nutzbare Bandbreite 
dagegen deutlich. Schon bei einem \textit{stride} von $2$ fällt die nutzbare Bandbreite
auf $169{,}15$ GB/s und für einen \textit{stride} von $4$ auf $101{,}76$ GB/s. 
Ab einem \textit{stride} von $8$ werden nur noch etwa $20$ Prozent des Peaks 
($\approx50$ GB/s) erreicht. Das liegt daran, dass mit zunehmendem \textit{stride} die 
Lanes eines Warps weiter auseinanderliegende Adressen lesen müssen. Dadurch müssen 
mehr Speichersegmente geladen werden, obwohl pro Segment nur ein kleiner Teil
tatsächlich genutzt wird.

\textit{Random Gather} hat mit $28{,}03$ GB/s die geringste nutzbare Bandbreite. 
Jede Lane liest praktisch eine zufällige Adresse, wodurch \textit{coalescing} und 
\textit{räumliche Lokalität} fast vollständig verloren gehen. Viele einzelne 
Speichertransaktionen und Cache-Misses senken dabei die effektive Bandbreite auf 
etwa $10$ Prozent des Peaks.

\section{\textit{Occupancy} und \textit{latency hiding}}

Im zweiten Experiment wurde die Blockgröße des \textit{coalesced} Kernels variiert. 
Für jede Blockgröße wurde die theoretische \textit{occupancy} mit
\texttt{cudaOccupancyMaxActiveBlocksPerMultiprocessor()} bestimmt:

\[
\mathrm{Occupancy}
= \frac{\mbox{aktive Warps}}{\mbox{maximal mögliche aktive Warps}}.
\]

Die Messwerte liegen in \texttt{outputs/occupancy.csv}; die Grafik ist in
Abbildung~\ref{fig:occupancy} dargestellt.

\begin{figure}[H]
  \centering
  \includegraphics[width=0.95\linewidth]{occupancy.png}
  \caption{Zusammenhang zwischen Blockgröße, \textit{Occupancy} und Bandbreite.}
  \label{fig:occupancy}
\end{figure}

\begin{table}[H]
  \centering
  \begin{tabular}{>{\raggedleft\arraybackslash}p{3.0cm}
                  >{\raggedleft\arraybackslash}p{3.0cm}
                  >{\raggedleft\arraybackslash}p{3.0cm}
                  >{\raggedleft\arraybackslash}p{2.5cm}}
    \hline
    \textbf{Blockgröße} & \textbf{Occupancy} & \textbf{GB/s} & \textbf{\% Peak} \\
    \hline
    32   & 0{,}500 & 183{,}24 & 67{,}36 \\
    64   & 1{,}000 & 251{,}30 & 92{,}38 \\
    128  & 1{,}000 & 251{,}32 & 92{,}39 \\
    256  & 1{,}000 & 251{,}33 & 92{,}39 \\
    512  & 1{,}000 & 251{,}30 & 92{,}38 \\
    1024 & 0{,}667 & 239{,}14 & 87{,}91 \\
    \hline
  \end{tabular}
  \caption{\textit{Occupancy-Sweep} für den \textit{coalesced} Kernel.}
  \label{tab:occupancy}
\end{table}

Bei einer Blockgröße von $32$ liegt die \textit{occupancy} nur bei $50$\% und die 
Bandbreite fällt auf $183{,}24$ GB/s. Das heißt, es sind zu wenige Warps in einem 
\textit{Streaming Multiprocessor (SM)}, um Speicher-Stalls vollständig zu überdecken. 
Ab einer Blockgröße von $64$ erreicht die  \textit{occupancy} stabil fast $100$\%. 
Die Bandbreite steigt dann auf $\approx251$ GB/s und bleibt bis zu einer Blockgröße von 
$512$ stabil. Ab einer Blockgröße von $1024$ sinkt die \textit{occupancy} wieder auf 
$\approx66,7$\%, da sich insgesamt weniger große Blöcke gleichzeitig auf einem 
\textit{Streaming Multiprocessor} befinden können; die
Bandbreite sinkt dadurch auf $239{,}14$ GB/s.

Das Experiment zeigt sog. \textit{latency hiding}: Da ein einzelner
Speicherzugriff mehrere Zyklen dauern kann, kann der \textit{Warp-Scheduler} bei einem 
\textit{stall} während der Verarbeitung eines Warps auf einen anderen Warp wechseln, 
solange sich genügend lauffähige Warps im \textit{Streaming Multiprocessor} befinden. Dadurch 
wird die Speicherlatenz eines einzelnen Warps zwar nicht vermieden, aber durch 
Parallelität überdeckt.

\subsection{GPU- und CPU-Strategien im Vergleich}

Die GPU vermeidet lange Speicherlatenzen durch einen hohen Durchsatz, indem sich viele 
Warps gleichzeitig in den \textit{SMs} befinden und Warp-Scheduling genutzt wird, um 
lange Speicherlatenzen zu maskieren.
Wie Experimente zu den Zugriffsmustern \textit{strided} und \textit{random gather} zeigen, 
sind GPU Caches im Vergleich zu CPU-Caches klein und können zufällige Zugriffe nicht in 
gleichem Umfang  abfedern. Deshalb entscheidet das Zugriffsmuster direkt über die 
erreichte Bandbreite.

Die CPU ist dagegen auf eine niedrige Latenz einzelner Instruktionsströme
optimiert. Sie nutzt hierfür mehrstufige Cache-Hierarchien, Hardware-Prefetching,
Out-of-Order-Execution und wenige starke Kerne. Sequenzielle Zugriffe
können dadurch sehr schnell sein, weil Daten rechtzeitig in den Cache geladen
werden. Falls zufällige Zugriffe Cache-Misses verursachen, sind diese auch für die CPU 
teuer.

\subsection{Folgen für das Datenlayout und GPU-Programme}

Aus den Messungen ergeben sich die folgenden Konsequenzen für effiziente
GPU-Programme:

\begin{itemize}
  \item Benachbarte Threads sollten benachbarte Speicheradressen lesen oder
        schreiben.
  \item Eine \textit{Structure-of-Arrays} ist für viele GPU-Zugriffsmuster günstiger als
        ein \textit{Array-of-Structures}, weil Felder dann zusammenhängend gelesen werden.
  \item \textit{Random Gather} und \textit{Scatter} sollten vermieden oder durch vorheriges
        Sortieren bzw. Umordnen der Daten reduziert werden.
  \item Daten sollten passend ausgerichtet und bei 2D-Strukturen gegebenenfalls
        gepaddet werden.
  \item Es müssen genug Threads bereitgestellt und passende Blockgrößen gewählt werden,
        damit sich immer genügend lauffähige Warps gleichzeitig in den SMs befinden können.
\end{itemize}

%\begin{verbatim}
%make
%make all
%make patterns
%make occupancy
%make stream
%make memory
%make report
%./build/bench mem 33554432 8 30
%./build/bench stream 67108864 20
%\end{verbatim}
